%% 03_methods.tex
\section{Methods}
\label{sec:methods}

\subsection{Dataset}
\label{sec:methods:dataset}

\paragraph{Base population.}
We generate 200 synthetic patients using Synthea~\citep{walonoski2018synthea}
with a fixed random seed, then apply a custom specialty-regimen overlay that assigns each
patient to one of three complexity tiers:
\begin{itemize}
  \item Tier 1 (Low): single long-acting injectable on a fixed recurring interval
    (e.g.\ every 8 weeks, modelling LAIs such as paliperidone palmitate~\citep{correll2024pp6m}).
  \item Tier 2 (Medium): cyclic regimen (on-for-$N$-weeks / off-for-$M$-weeks)
    requiring cycle-position reasoning.
  \item Tier 3 (High): multi-drug regimen (biologic primary + oral adjunct + PRN
    rescue) with staggered or interdependent schedules.
\end{itemize}
The final distribution is 63 Tier-1, 63 Tier-2, and 74 Tier-3 patients.  All 200 FHIR bundles
are validated against a reference R4B pre-write conformance validator
(\path{mh_integration/r4b_validator.py}) using \texttt{fhir.resources} $\geq$8.0; all
200 bundles pass with no mandatory element failures.  Only Tier-3 patients have an oral
adjunct or a PRN rescue component; this fact drives the N/A structure of the question bank
described below.

\paragraph{Paired representations.}
Each patient is represented in two parallel forms derived from the same FHIR bundle:
\begin{enumerate}
  \item FHIR bundle (structured): Synthea base plus the specialty overlay
    (MedicationRequest, MedicationAdministration, Medication, CarePlan, Encounter, Patient,
    Practitioner), frozen at release tag \texttt{dataset-freeze-v1}
    (manifest SHA-256 \texttt{4d0da93e\ldots}).  The structured systems index the
    \emph{entire} bundle, including the Synthea base record, not only the overlay
    resources.
  \item LLM narrative: generated from the regimen overlay by Llama~3.3-70B via Groq with
    a fixed system prompt (hash \texttt{518bc71d87b7}) and temperature 0.0.  The prompt
    instructs the generator to state, per component, the first three and last three
    administration dates, one mid-window date, and the \emph{total administration count}.
    A deterministic templated narrative generated from the same overlay fields, with no
    LLM, serves as a second narrative source.
\end{enumerate}

\paragraph{What the pairing does and does not hold constant.}
The two representations are derived from one source, so no fact in the narrative lacks a
FHIR counterpart.  They are not, however, information-equivalent, and we no longer describe
the design as holding information content constant.  Three asymmetries matter for
interpretation.  First, the narrative is a curated summary of eight prompted entity
classes over the overlay only, whereas the bundle also contains the full Synthea base
record (encounters, conditions, observations) that is irrelevant to every question.
Second, the narrative states a precomputed total administration count per component; the
FHIR bundle contains the individual MedicationAdministration resources but no count.
Third, the narrative omits all but seven administration dates per component, while the
bundle contains every one (a daily oral component accumulates over 250 events).  LLM
narratives have a median length of 84 words (range 22--411; templated median 108) and
therefore fit inside a single 500-token retrieval chunk, so the narrative arm's ``top-5
retrieval'' returns the whole document and performs no selection.

\paragraph{Fidelity audit.}
Narrative quality is verified by a programmatic audit that recovers eight prompted entity
classes (medication name, regimen start/end, tier integer, tier description,
administration dates, administration count, dosage schedule) directly from FHIR
resources.  The audit verifies that entities explicitly listed in the generation prompt
were preserved in the narrative; it does not claim recovery of FHIR resources outside the
prompted set.  After one prompt-iteration step to add mid-window administration dates,
both generators reached 100\% weighted prompted-entity recall across all 200 patients and
all 2{,}403 checks.

\paragraph{Question bank.}
We generate 13,800 questions (200 patients $\times$ 69 templates) across five PSP
adherence-indicator families.  Each template carries one canonical question string and
two stored paraphrases; \emph{only the canonical string was evaluated}, so the 13,800
rows are 69 logical questions per patient with no paraphrase duplication.  Tables and
figures use the technical family names (in monospace) for grep-ability against the
released code; the PSP-clinical glosses follow each name in parentheses.
\begin{enumerate}
  \item \texttt{temporal\_lookup} (next-dose lookup; $n=3{,}200$):
    ``What is the next scheduled administration on the reference date?''
  \item \texttt{regimen\_aggregation} (dose-history aggregation, MPR numerator;
    $n=2{,}400$): ``How many doses have been administered in the past 90 days?''
  \item \texttt{regimen\_compliance} (coverage-window reasoning, PDC;
    $n=4{,}200$): ``What proportion of the past 90 days was the patient covered?''
  \item \texttt{temporal\_comparison} (missed-dose detection; $n=2{,}400$):
    ``Is the most recent administration before, on, or after the reference date?''
  \item \texttt{cross\_resource} (persistence and component-presence questions;
    $n=1{,}600$): ``Does the regimen include an oral adjunct component?''
\end{enumerate}
All ground-truth values are computed programmatically from FHIR bundles using pure-Python
adherence metric functions (MPR, PDC, is\_persistent, consecutive\_missed\_doses); no manual
annotation and no LLM-as-judge is used anywhere in evaluation.  The question bank SHA-256
is \texttt{2216f58e\ldots}.%
\footnote{Every template fires for every patient, so $13{,}800 = 200 \times 69$ exactly.
The tier-by-tier counts (Tier~1: 4{,}347; Tier~2: 4{,}347; Tier~3: 5{,}106) reflect only
the patient distribution (63 / 63 / 74).  No template is tier-conditional.}

\paragraph{N/A structure of the question bank.}
Because every template fires for every patient and 37 of the 69 templates address the
adjunct or rescue component, the correct answer is N/A whenever the component does not
exist.  Table~\ref{tab:na_prevalence} gives the resulting prevalence.  Overall 5{,}623 of
13{,}800 ground truths (40.7\%) are N/A; at Tiers~1 and~2 the majority of questions have
N/A as the correct answer, whereas at Tier~3 only 13.8\% do.  A constant predictor that
always answers N/A therefore scores 54.2\%, 58.9\%, and 13.8\% at Tiers~1--3 and 40.7\%
overall.  All accuracy results in \S\ref{sec:results} are reported both in aggregate and
stratified by this variable.

\begin{table}[htbp]
\centering
\caption{Prevalence of N/A ground truth by tier.  ``Always N/A'' is the accuracy of a
constant predictor that returns N/A for every question.}
\label{tab:na_prevalence}
\begin{tabular}{lrrrr}
\toprule
Tier & Questions & N/A ground truth & Substantive ground truth & Always-N/A accuracy \\
\midrule
1 & 4{,}347 & 2{,}357 & 1{,}990 & 54.2\% \\
2 & 4{,}347 & 2{,}562 & 1{,}785 & 58.9\% \\
3 & 5{,}106 &   704   & 4{,}402 & 13.8\% \\
\midrule
All & 13{,}800 & 5{,}623 & 8{,}177 & 40.7\% \\
\bottomrule
\end{tabular}
\end{table}

\paragraph{Reference date.}
Each question row carries a per-patient \texttt{reference\_date} field, and 7{,}800 of the
13{,}800 canonical question strings refer to ``the reference date''.  The harness passes
only the question string and the patient identifier to each system; the reference date
is not interpolated into the question, not included in the answer prompt, and not present
in the narratives (it appears by coincidence in one of 200 LLM narratives).  We discovered
this during the post-review audit.  Its consequences are quantified in
\S\ref{sec:results:flaws}.

\subsection{Retrieval Systems}
\label{sec:methods:systems}

All systems share controlled variables (Table~\ref{tab:controlled}):
embedding model \texttt{BAAI/bge-large-en-v1.5} (local via sentence-transformers),
answer LLM \texttt{qwen/qwen3-32b} via Groq (temperature 0.0, single run per question),
ChromaDB local vector store, $k=5$ retrieved chunks, chunk size 500 tokens, overlap 50 tokens.

\begin{table}[htbp]
\centering
\caption{Controlled variables held constant across Systems A, B, C, A-T and N.}
\label{tab:controlled}
\begin{tabular}{ll}
\toprule
Parameter & Value \\
\midrule
Embedding model & \texttt{BAAI/bge-large-en-v1.5} (local) \\
Answer LLM & \texttt{qwen/qwen3-32b} via Groq Developer plan \\
Reasoning mode & Groq default for this model: thinking enabled, \texttt{<think>} block returned inline \\
Temperature & 0.0 \\
Runs per question & 1 \\
Top-$k$ & 5 \\
Chunk size & 500 estimated tokens (whitespace words $\times$ 1.3) \\
Chunk overlap & 50 estimated tokens \\
Vector store & ChromaDB (local) \\
Answer max tokens & 512 (thinking and answer share the budget) \\
\bottomrule
\end{tabular}
\end{table}

\paragraph{Answer LLM configuration.}
Qwen3-32B is a hybrid reasoning model.  The Groq request set only \texttt{model},
\texttt{messages}, \texttt{temperature} and \texttt{max\_tokens}; no
\texttt{reasoning\_effort} or \texttt{reasoning\_format} parameter was passed, so the
model ran with thinking enabled and returned its reasoning inline between
\texttt{<think>} tags, sharing the 512-token output budget with the final answer.  Every
one of the 69,000 responses across all five arms begins with a \texttt{<think>} block.
This was not a deliberate design choice and we treat it as a pipeline defect
(\S\ref{sec:results:flaws}).

\paragraph{System A (Narrative RAG).}
LLM narratives are split into 500-token chunks with 50-token overlap, embedded, and indexed.
At query time the top-5 chunks are retrieved by cosine similarity and concatenated into the
answer prompt.  Because every narrative is a single chunk, the context is always the
full narrative (mean 490 prompt tokens, maximum 817).

\paragraph{System B (Structured RAG, Naive).}
Each FHIR resource in the bundle is serialised to indented JSON, chunked with the shared
chunker, embedded, and indexed.  Retrieval is identical to System~A.  Mean prompt
tokens per call are 4,463 (median 2,597; maximum 128,385).  The mean exceeds five nominal
500-token chunks because the chunker estimates tokens as whitespace words $\times$ 1.3
while Groq bills with the Qwen3 tokenizer, and indented JSON has far more tokenizer tokens
per whitespace word than prose; the maximum arises from resources with no sentence
boundaries that the sentence-based chunker cannot split.

\paragraph{System C (Structured RAG, Resource-Aware).}
Two mechanisms are added to System~B, and nothing else.  (1)~A deterministic regex router
maps the \emph{question string} (not the question bank's family label) to a prioritised
resource-type allowlist that is applied as a ChromaDB metadata filter before similarity
ranking; the last router entry is a catch-all that applies no filter, and the system
falls back to unfiltered retrieval when the filtered pool returns fewer than $k$ hits.
(2)~After top-$k$ retrieval, the reference graph built at index time is walked one hop
(e.g.\ MedicationRequest $\rightarrow$ Medication, Practitioner, Patient) and up to ten
expansion chunks are appended in a marked section.  A temporal pre-filter on resource
dates is also applied where the router selects date-bearing resource types.  Mean prompt
tokens: 2,175 (median 1,649).  Because the router sees only the question text, System~C
receives no oracle information about the question family.

\paragraph{Additional arms.}
System~N (no-retrieval baseline) uses the same answer LLM, system prompt, temperature,
and max-tokens as A/B/C but receives an empty retrieved context.  System~A-T
(templated-narrative arm) uses the System~A pipeline verbatim but indexes the
deterministic templated narratives instead of the LLM narratives.  Both narrative sources
pass the same prompted-entity fidelity audit at 100\% recall.

\subsection{Evaluation}
\label{sec:methods:eval}

\paragraph{Accuracy.}
The scorer (\texttt{eval/score.py}) applies deterministic type-detection rules to classify
each ground truth as date, numeric, boolean/N-A, list, or free-text, then applies
type-specific scoring.  For date and numeric ground truths the scorer searches the
\emph{entire} response text, including the \texttt{<think>} block, and takes the first
ISO date or first number it finds (2\% relative tolerance for floats).  Boolean and N/A
ground truths are matched canonically.  Free-text ground truths (component names such as
``oral adjunct'', comparators such as ``same-day'' or ``before'') are scored by checking
that a set of bundle-derived expected entities all appear in the response, not by matching
the ground-truth string itself.  The primary metric is exact-match rate.  The reach of the
scorer into the reasoning trace, and the behaviour of the free-text path, are audited in
\S\ref{sec:results:flaws}.

\paragraph{Retrieval recall@k.}
For Systems~B and~C only (System~A's narrative chunks carry no FHIR resource IDs), recall@$k$
is computed by checking whether the top-$k$ retrieved chunk IDs include the patient's primary
MedicationRequest or specialty CarePlan (a conservative heuristic; see Limitations).  Because
only $k=5$ chunks are retrieved, recall@10 is identical to recall@5 by construction and is
not reported.

\paragraph{Statistical inference.}
The 13,800 questions are not independent: each of the 200 patients contributes 69
questions that share a regimen, a tier, and a narrative.  Our primary inference is
therefore a \emph{patient-level cluster bootstrap}: 10,000 resamples of the 200 patients
with replacement (seed 42), recomputing each paired accuracy difference on the resampled
question set, with percentile 95\% CIs.  We also report the question-level paired
bootstrap (10,000 resamples, paired by question id) for comparison with the previous
version of this paper, and McNemar's uncorrected $\chi^2$ on the discordant pairs as a
descriptive statistic; McNemar assumes independent pairs, so we do not rely on it for
inference.  Wilson score 95\% CIs are reported for per-cell proportions.  Three
pre-specified comparisons (A vs.\ B, A vs.\ C, B vs.\ C); no multiple-comparison
correction.  All N/A-stratified analyses are post hoc and are labelled as such.

\paragraph{Output-budget audit.}
Truncation is measured directly from the billed \texttt{completion\_tokens} field
(a response is truncated if it used the full 512-token budget) and from whether the
response contains a closing \texttt{</think>} tag.  A regex-based error taxonomy over a
50-failure sample per system was reported in the previous version of this paper; it was
computed on a CSV in which responses had been cut to 500 characters for readability, so
its ``reasoning-truncated'' category conflated CSV truncation with generation truncation.
We retain the script in the repository but no longer report the taxonomy as a result.

\paragraph{Exploratory feature-extraction arm.}
Three patient-level feature sets are constructed for a synthetic adherence label.
FS-Structured consists of 19 FHIR-derived numeric features (MPR-90d/180d/full,
PDC-90d/full, gap statistics, event counts, days-since-last-dose, tier one-hot,
resource-type counts, age).  FS-Narrative consists of 6 heuristic regex features on LLM
narratives.  FS-Aware consists of 8 per-patient aggregates from System~C retrieval traces.
LightGBM ($n\_\text{estimators}=200$, learning rate 0.05, class-weight balanced) and
logistic regression ($C=1.0$, class-weight balanced) are evaluated via stratified 5-fold
cross-validation, with patient-level paired bootstrap AUC CIs from out-of-fold predictions.
The three feature sets differ in engineering effort as well as in source, so this arm
compares pipelines, not representations; it is reported as exploratory in
\S\ref{sec:results:o6}.

\paragraph{API and cost disclosure.}
Narrative generation (200 calls) used the Groq free tier.  The full evaluation matrix
($\approx$69{,}000 calls across five arms plus the 1,485-call budget sweep) used the Groq
Developer paid plan.  Observed API cost for the three primary arms was \$32.68~USD
(System~A: \$4.74; System~B: \$21.07; System~C: \$6.88).  All calls used
temperature~0.0; responses are cached and reproducible.  No PHI was transmitted: the
dataset is fully synthetic.
